\section{Clifford点における低階数切断の局所被覆}
\label{sec:clifford-rank-slice}
この節の基点は式\eqref{eq:clifford-rules}の$C_0$であり、
全Cliffordチャートや全低階数領域を一度に分類するものではない。

\begin{lemma}[有限双対のArtin平坦性から元代数の平坦性]\label{lem:dual-artin-flatness}
局所Artin環$B$上の二次四窓について、双対$E_B$が各次数で有限自由であり、
特殊ファイバー$E_0$がKoszulなら、$E_B$は相対Koszulである。
その二次双対の各成分も自由で、特殊ファイバーと同じHilbert元数を持つ。
特に$C_0$を特殊ファイバーとする有限Frobenius候補のArtin変形は、
元代数の成分ごとに平坦な変形を与える。
\end{lemma}
\begin{proof}
補題\ref{lem:relative-four}と同じGaussian消去を、今回は$E_B$のbar複体に適用する。
内部次数を固定すれば有限自由複体であり、剰余体上の可縮対を単元で消去すると、
ホモロジー次数と内部次数が等しい一つの次数だけに自由加群が残る。
従って対角Torは自由で、それを双対にした二次双対の各成分も自由となる。
有限Frobenius候補の定数階数条件は次数$0,\ldots,4$での自由性を保証する。
次数5の特殊ファイバーは零だからNakayama補題で次数5自体が零となり、全高次も零である。
従ってこの場合も全次数で前半の議論を適用できる。
\end{proof}

$w_0$を$C_0$の最上traceテンソルとし、$0123$の係数を1に正規化する。
第三の二対二行列を$M_0$とする。行番号は$4a+c$、列番号は$4b+d$である。
$M_0$は対称な階数6の$16\times16$行列であり、通常の有理数nullspace算法で
$K\in\operatorname{Mat}_{16,10}$を$\ker M_0$の列基底として選ぶ。
一次変形のtrace変化を$\dot w$、対応する行列を$\dot M$とすると、
7次小行列式の一次条件は$K^{\mathsf T}\dot M K=0$である。

\begin{proposition}[低階数条件の一次切断と障害のPfaffian]\label{prop:clifford-rank-derivative}
既定の偶基底$e_0,\ldots,e_8$と交代基底$o_0,\ldots,o_5$に対し、
写像$\dot w\mapsto K^{\mathsf T}\dot M K$は全偶方向と全基底変更方向で零となり、
交代6次元上では単射である。
従って四周期平坦表示の75次元接空間は、低階数条件によって69次元に切られる。
交代方向の六つの$10\times10$行列は共通の6次元核を持ち、商の基底を適切に取ると
\[
 S(o)=\begin{pmatrix}
 0&o_5&o_2-o_4-o_5&-o_4\\
 -o_5&0&o_1&o_1-o_3\\
 -o_2+o_4+o_5&-o_1&0&-o_0-o_1+o_3\\
 o_4&-o_1+o_3&o_0+o_1-o_3&0
 \end{pmatrix}
\]
となる。このPfaffianは既知の二次障害$Q$の負である。
\end{proposition}
\begin{proof}
双対関係を順序付き二次語$m$で番号付けすると、その基底は
$m+\sum_r F_r(m)\,\operatorname{desc}(r)$である。
各一次代表$\alpha$について、この基底の変化を三つの位置へ挿入した行列を$D_1$、
基点の挿入行列を$D_0$とすれば、trace変化は
\[
 D_0^{\mathsf T}\dot w=-D_1^{\mathsf T}w_0,\qquad \dot w_{0123}=0
\]
で一意に決まる。ここで$D_0$は$256\times480$行列である。
480本の式と正規化を有理数上で解き、各解を元の全式へ代入して確認する。
\path{computations/verify_clifford_rank_slice.py}はこの481本の式を15代表について一括で検証する。

同コードで偶9方向の像は零、交代6方向の像は階数6となる。
100成分を行に並べた交代微分行列では、0始まりの行番号
$12,15,19,25,29,59$の小行列式が$-1$である。
64個の初等基底変更方向についても像が零であることを直接確認する。
定理\ref{thm:clifford-tangent}の複体より表示の接次元は$60+9+6=75$なので、切断後は69となる。

共通核を除くと、有理数rrefで選んだ$K$座標の番号$1,2,5,9$が商の基底を与え、
表示した$S(o)$を得る。最後に
\[
 \operatorname{Pf}S=S_{01}S_{23}-S_{02}S_{13}+S_{03}S_{12}
 =-o_0o_5-o_1o_2+o_2o_3-o_3o_4=-Q
\]
を展開すればよい。これは既知の障害の行列的解釈であり、
高次障害や全変形環の新しい決定を主張するものではない。
\end{proof}

\begin{theorem}[Clifford点の低階数近傍は対称チャートで被覆される]\label{thm:clifford-low-rank-local}
テンソル復元定理の有限Frobenius候補空間に、第三の7次小行列式の消滅を課した
局所有限型スキームを$\mathfrak T_{\leq6}$とする。
このスキームは$[w_0]$で滑らかな次元69の表示空間であり、
$[w_0]$を含むある$G$不変Zariski開部分は対称Frobeniusチャートで被覆される。
従ってこの開部分の候補は全てKoszul AS正則であり、
\texttt{LOW-UNCLASSIFIED}はこの近傍に現れない。

より具体的に、既定の9個の偶一次代表をそのまま基点の関係へ加えた9変数族は、
基点近傍で全ての基点付き低階数Artin変形を、添字保持同型を除いて含む。
これは低階数切断内の局所主張であり、無制限のClifford変形環や全低階数領域の被覆ではない。
\end{theorem}
\begin{proof}
定数階数の行列条件は、選んだ非零小行列式の近傍でkernelとcokernelを有限自由にする。
従ってテンソル復元条件は局所有限型の階数層として定義でき、
四窓への写像と最上交差空間の生成線への逆写像はこの階数層上で正則である。
Artin環上でも最上交差空間は自由階数1なので、この写像は無限小にも単射である。
補題\ref{lem:dual-artin-flatness}により、その全Artin接方向は平坦PBW四窓の接方向に含まれる。
前命題から$\dim T_{[w_0]}\mathfrak T_{\leq6}\leq69$となる。

一方、偶基底の9代表は全て二次平方への係数変更だけである。
従ってこの9変数族の関係空間は対称テンソル内にとどまり、双対は可換な四二次式の商である。
基点で完全交叉なので、その非空開近傍$V\subset\mathbb A^9$は
定理\ref{thm:symmetric-frobenius-chart}のチャートへ入る。
特に$V$全体が平坦かつ階数6の族である。
滑らかな源$G\times V$から$\mathfrak T_{\leq6}$への射を考えると、
基底変更の60方向と独立な偶9方向により接写像の階数は69である。
上の上界と合わせて接空間を全射する。
補題\ref{lem:smooth-source}より、標的はこの点で滑らかな次元69となる。

両空間が滑らかで接写像が全射なので、射は対応点で滑らかである。
滑らかな射の像は開であり、その像は全て対称チャートの像に含まれる。
この開近傍を$G$で飽和すれば主張する不変開部分を得る。
形式的な局所完全性も補題\ref{lem:smooth-source}による。
無限小の裸の同型と標準四周期表示の比較は定理\ref{thm:clifford-tangent}と
補題\ref{lem:relative-four}で処理済みであり、全局所対象の漏れはない。
\end{proof}
\begin{remark}
69は基底付き表示空間の局所次元で、stack次元や全Clifford成分の次元ではない。
基点付き同型を除いた接方向は9であるが、有限自己同型で割る粗モジュライの滑らかさは主張しない。
また、低階数の閉条件を課さない全Clifford点には交代変形と障害$Q$が残る。
\end{remark}
